% Requires: \usetikzlibrary{positioning,arrows.meta}
\begin{figure}[H]
  \centering
  \resizebox{0.99\textwidth}{!}{%
  \begin{tikzpicture}[
      >={Stealth[length=1.8mm]},
      res/.style={draw=black!70,line width=0.4pt,fill=black!6,
        rounded corners=2pt,align=center,inner sep=3.5pt,
        font=\footnotesize},
      dep/.style={draw=black,line width=0.5pt,->}
    ]
    % row 1: the structural toolkit and the two reduction rules
    \node[res,text width=2.6cm] (FM) at (2.6,0)
      {feasibility and monotonicity\\
       \textcolor{black!60}{\scriptsize Props.~\ref{prop:feas}
       and~\ref{prop:mono}}};
    \node[res,text width=2.0cm] (OR) at (7.6,0)
      {routing oracle\\
       \textcolor{black!60}{\scriptsize Prop.~\ref{prop:oracle}}};
    \node[res,text width=2.4cm] (AGG) at (12.4,0)
      {aggregation and capping\\
       \textcolor{black!60}{\scriptsize Prop.~\ref{prop:aggregation},
       Obs.~\ref{obs:capping}}};
    % row 2: constructions, algorithms, certificates, compression
    \node[res,text width=2.8cm] (CC) at (1.5,-2.1)
      {the hardness constructions\\
       \textcolor{black!60}{\scriptsize Thms.~\ref{thm:cover},
       \ref{thm:products}, \ref{thm:partition}, \ref{thm:numeric},
       Cors.~\ref{cor:mirror}, \ref{cor:partitionmirror}}};
    \node[res,text width=2.0cm] (FPT) at (4.6,-2.1)
      {design enumeration\\
       \textcolor{black!60}{\scriptsize Thm.~\ref{thm:fpt}}};
    \node[res,text width=2.2cm] (XP) at (7.2,-2.1)
      {\XP\ branch-and-bound\\
       \textcolor{black!60}{\scriptsize Thm.~\ref{thm:xp},
       Alg.~\ref{alg:xpbb}}};
    \node[res,text width=2.0cm] (BC) at (9.8,-2.1)
      {Benders certificates\\
       \textcolor{black!60}{\scriptsize Prop.~\ref{prop:duals}}};
    \node[res,text width=2.4cm] (CP) at (12.4,-2.1)
      {polynomial compression\\
       \textcolor{black!60}{\scriptsize Prop.~\ref{prop:ftcompress}}};
    % row 3: the classification
    \node[res,text width=2.4cm] (W2) at (0.9,-4.2)
      {$\mathrm{W}[2]$ and ETH bounds\\
       \textcolor{black!60}{\scriptsize Thms.~\ref{thm:w2}
       and~\ref{thm:eth}}};
    \node[res,text width=2.2cm] (NK) at (3.5,-4.2)
      {no polynomial kernel\\
       \textcolor{black!60}{\scriptsize Thm.~\ref{thm:nokernel}}};
    \node[res,text width=1.9cm] (DI) at (6.1,-4.2)
      {size dichotomy\\
       \textcolor{black!60}{\scriptsize Thm.~\ref{thm:dichotomy}}};
    % dependencies
    \draw[dep] (FM.south) -- (CC.north);
    \draw[dep] (FM.south) -- (FPT.north);
    \draw[dep] (FM.south) -- (XP.north);
    \draw[dep] (OR.south) -- (FPT.north);
    \draw[dep] (OR.south) -- (XP.north);
    \draw[dep] (OR.south) -- (BC.north);
    \draw[dep] (OR.south) -- (CP.north);
    \draw[dep] (AGG.south) -- (CP.north);
    \draw[dep] (CC.south) -- (W2.north);
    \draw[dep] (CC.south) -- (NK.north);
    \draw[dep] (CC.south) -- (DI.north);
    \draw[dep] (FPT.south) -- (DI.north);
  \end{tikzpicture}%
  }
  \caption{The results of the paper as a diagram, each box naming its
  numbered statements. A solid arrow points from a result to one whose
  proof uses it, and no other relation is drawn. The structural layer
  feeds every branch, the hardness constructions carry the whole
  negative side, from the $\mathrm{W}[2]$ and ETH bounds
  (Theorems~\ref{thm:w2} and~\ref{thm:eth}) through the kernel
  exclusion (Theorem~\ref{thm:nokernel}) to the negative half of the
  dichotomy (Theorem~\ref{thm:dichotomy}), and the oracle feeds the
  algorithms, the cuts and the compression
  (Prop.~\ref{prop:ftcompress}).}
  \label{fig:resultsmap}
\end{figure}
